\documentclass[10pt,a4paper]{amsart}
\usepackage[margin=19mm]{geometry}
\usepackage{amsmath,amssymb,mathtools}
\usepackage[scr=rsfs,cal=euler]{mathalfa}
\usepackage{xfrac}
\usepackage[unicode,hidelinks,bookmarks=false]{hyperref}
\newcommand{\ie}{i.e.\ }
\newcommand{\cf}{cf.\ }
\newcommand{\resp}{respectively\ }
\newcommand{\Subsection}{\subsection}
\providecommand{\nobreakdash}{\nobreak\hbox{-}}
\setlength{\emergencystretch}{2em}
\multlinegap0pt
\usepackage{amssymb}
\usepackage{float}
\usepackage{enumerate}
\usepackage{tikz}
\newtheorem{prop}{Proposition}
\newtheorem{claim}{Claim}
\usepackage{cleveref}
\crefname{prop}{Proposition}{Propositions}
\crefname{claim}{Claim}{Claims}
\title{On the PPW Conjecture for Hopf-Symmetric Sets in Non-compact Rank One Symmetric Space}
\author{Yusen Xia}
\date{Source: arXiv:2505.18439v3 (CC BY 4.0)}
\begin{document}
\maketitle

\noindent\emph{Source and licence.} On the PPW Conjecture for Hopf-Symmetric Sets in Non-compact Rank One Symmetric Space. Version arXiv:2505.18439v3 (CC BY 4.0), distributed by arXiv under the Creative Commons Attribution 4.0 International licence, http://creativecommons.org/licenses/by/4.0/, as recorded in the arXiv licence metadata for this exact version and shown on the abstract page https://arxiv.org/abs/2505.18439v3. Full licence notice: \texttt{licenses/arxiv-2505.18439-CC-BY-4.0.txt}.\par\smallskip

\noindent\textit{Editorial scope.} Counted unit: the article's prop g1decreasinglogconcave (source0.tex lines 314--316). The article's complete original proof of it is delivered with it (source0.tex lines 317--351, one \texttt{proof} environment of 35 lines). No mathematics is written, repaired or re-derived: the statement and the proof are the article's own lines, in the article's own notation, and no retained line is substituted or re-flowed.  Retained uncounted as the source-local context the proof needs: lines 310--312.  Nothing was omitted from any retained span, so the package carries no omission record.  No retained line contains a citation, so the wrapper carries no bibliography and no bibliography entry.  No retained line contains a figure environment, an image-inclusion command or drawing code, so no image is delivered and the package declares no asset dependency.  Editorial formatting only, in addition: an AMS \texttt{amsart} preamble that restates the article's own declarations and macros used by the retained lines, namely the environments \texttt{prop}, \texttt{claim} on separate counters, together with the standard packages those lines need.  The retained environments print in this document as Proposition 1, Claim 1 in order of appearance, whereas the article prints the same items with the numbering of its own full document.  The complete native source of the article as distributed in the arXiv e-print for this exact version is bundled unchanged as \texttt{sources/178695/source0.tex}.
\begin{equation}\label{FirstEigenODE}
    g_1^{\prime \prime}+H(r) g_1^{\prime}+\lambda_1(B)  g_1=0 \textit{ with } g_1'(0) = 0 \textit{ and } g_1(R) = 0.
\end{equation}

\begin{prop}\label{g1decreasinglogconcave}
    The first Dirichlet eigenfunction of a geodesic ball is decreasing and log-concave.
\end{prop}

\begin{proof}
    We only prove the statement for balls in noncompact spaces here; for compact spaces, just replace the mean curvature function $H$ appropriately. 
    \par As $H(r) = \frac{J'(r)}{J(r)}$, multiplying $J(r)$ on both sides of \cref{FirstEigenODE} yieds
    \[\left( J g_1'\right)' = -\lambda_1 Jg_1<0.\] Hence $Jg_1'$ is decreasing in $r$. $Jg_1'\leq \lim_{r\to 0^+} J(r)g_1(r)= 0$, thus $g'_1 \leq 0$.
    \par For log-concavity, let $\varphi=\left(\log g_1\right)^{\prime}=\frac{g_1^{\prime}}{g_1}$ then $\varphi(0)=0$, $\varphi<0$ on $(0, R)$. We compute that
    \[\begin{aligned}
\varphi^{\prime} & =\frac{g_1^{\prime \prime}}{g_1}-\left(\frac{g_1^{\prime}}{g_1}\right)^2\\
&=-H(r)\left(\frac{g_1^{\prime}}{g_1}\right)-\lambda_1(B)-\left(\frac{g_1^{\prime}}{g_1}\right)^2 \\
& =-\left(\frac{kn-1}{\tanh r}+(k-1) \tanh r\right) \varphi-\lambda_1(B)-\varphi^2
\end{aligned}\] At $r=0:$
\[\begin{aligned}
     \varphi^{\prime}(0) &=-\lambda_1(B)-\lim _{r \rightarrow 0^{+}}\left(\frac{kn-1}{\tanh r}+(k-1) \tanh r\right) \varphi(r)\\
    &=-\lambda_1(B)-(k n-1) \varphi^{\prime}(0) .
\end{aligned}\]
Hence $\varphi'(0)<0.$ 
\begin{claim}
  $\varphi^{\prime}<0$ on $[0, R).$  
\end{claim}
Otherwise there is some $r_1\in (0,R)$ such that
\begin{equation*}
\left\{\begin{array}{l}
\varphi^{\prime}(r)<0 \text { on } (0,r_1)\\
\varphi^{\prime}\left(r_1\right)=0 . \\
\varphi^{\prime \prime}\left(r_1\right) \geqslant 0
\end{array}\right.
\end{equation*}
But \[\varphi^{\prime}=-H(r) \varphi-\lambda_1(B)-\varphi^2\] implies that
\[\begin{aligned}
\varphi^{\prime \prime} & =-H^{\prime}(r) \varphi-H(r) \varphi^{\prime}-2 \varphi \varphi^{\prime} \\
& =\lambda_1\left(S_r\right) \varphi-H(r) \varphi^{\prime}-2 \varphi \varphi^{\prime} .
\end{aligned}\]
At $r=r_1$: \[0 \leq \varphi^{\prime \prime}\left(r_1\right)=\lambda_1\left(S_{r_1}\right) \varphi\left(r_1\right)<0.\] 
A contradiction. This completes the proof of the claim and the lemma.

\end{proof}

\end{document}
